%%%%%%%%%%%%%%%%%%%%%%%%%%%%%%%%%%%%%%%%%%%%%%%%%%%%%%%%%%%%%
%\sectioncontents
%%%%%%%%%%%%%%%%%%%%%%%%%%%%%%%%%%%%%%%%%%%%%%%%%%%%%%%%%%%%
%%%%%%%%%%%%%%%%%%%%%%%%%%%%%%%%%%%%%%%%%%%%%%%%%%%%%%%%%%%%
\begin{slide}
\section{REPRESENTATION PHYSIQUE DANS ORACLE}
\end{slide}
%%%%%%%%%%%%%%%%%%%%%%%%%%%%%%%%%%%%%%%%%%%%%%%%%%%%%%%%%%%%
%%%%%%%%%%%%%%%%%%%%%%%%%%%%%%%%%%%%%%%%%%%%%%%%%%%%%%%%%%%%
\begin{slide}
\subsection{Repr\'esentation physique dans ORACLE V7}

Les principales structures physiques utilis\'ees dans ORACLE sont\,:
\begin{enumerate}
  \item Le {\bf bloc} est l'unit\'e physique d'E/S.
     La taille d'un bloc ORACLE est un
        multiple de la taille des blocs du syst\`eme sous-jacent.
  \item L'{\bf extension} est un ensemble de blocs contigus contenant
        un m\^eme type d'information.
  \item Le {\bf segment} est un ensemble d'extensions stockant un 
        objet logique (une table, un index ...). 
\end{enumerate}
\end{slide}
\begin{slide}
\subsection{Tables, segments, extensions et blocs}

\mywidthfigure{8cm}{oraphys.eps}

\end{slide}
\begin{slide}
\subsection{Le segment ORACLE}

Le segment est la zone physique contenant un objet logique. Il
existe quatre types de segments\,:
\begin{enumerate}
  \item Le segment de donn\'ees (pour une table ou un $cluster$).
  \item Le segment d'index.
  \item Le $rollback\ segment$ utilis\'e pour les transactions.
  \item Le segment temporaire (utilis\'e pour les tris par exemple).
\end{enumerate}
\end{slide}
\begin{slide}
\subsection{Base ORACLE, fichiers et  $TABLESPACE$}

\begin{enumerate}
 \item {\bf Physiquement}, une base ORACLE est un ensemble de fichiers.
 \item {\bf Logiquement}, une base est divis\'ee par l'administrateur
   en $tablespace$. Chaque $tablespace$ consiste en un (au moins)
        ou plusieurs fichiers.
\end{enumerate}
La notion de $tablespace$ permet\,:
\begin{enumerate}
  \item De contr\^oler l'emplacement physique des donn\'ees.
        (par ex.\,: le dictionnaire sur un disque,
        les donn\'ees utilisateur sur un autre).
  \item de faciliter la gestion (sauvegarde, protection, etc).
\end{enumerate}
\end{slide}
\begin{slide}

\myfigure{oratbspace.eps}

\end{slide}
\begin{slide}
\subsection{Stockage des donn\'ees}

Il existe deux mani\`eres de stocker une table\,:
\begin{enumerate}
  \item {\bf Ind\'ependamment}. Un segment est alors automatiquement
        allou\'e \`a la table. Il est possible de sp\'ecifier
        des param\`etres pour ce segment\,:
       \begin{enumerate}
             \item Sa taille initiale
             \item Le pourcentage d'espace libre dans chaque bloc.
             \item La taille des extensions.
       \end{enumerate}
  \item {\bf Dans un} $cluster$.
 \end{enumerate}
% Parler du ROWID
\end{slide}
\begin{slide}
\subsection{Stockage des n-uplets}

En r\`egle g\'en\'erale un n-uplet est stock\'e dans un seul bloc. 
L'adresse physique d'un n-uplet est le $ROWID$ qui se d\'ecompose
en trois parties\,:
\begin{enumerate}
   \item Le num\'ero du n-uplet dans la page.
   \item Le num\'ero de la page, relatif au {\bf fichier} dans lequel
         se trouve le n-uplet.
   \item Le num\'ero du fichier.
\end{enumerate}
Exemple\,: 00000DD5.000.001 est l'adresse du premier n-uplet
du bloc DD5 dans le premier fichier.
\end{slide}
\begin{slide}
\subsection{Structures de donn\'ees pour l'optimisation}

ORACLE 7 propose trois structures pour l'optimisation de requ\^etes\,:
\begin{enumerate}
  \item Les index
  \item Les ``regroupements'' de tables ($Cluster$).
  \item Le hachage.
\end{enumerate}
\end{slide}
\begin{slide}
\subsection{Les index ORACLE}

On peut cr\'eer des index sur tout attribut (ou tout ensemble
d'attributs) d'une table. ORACLE utilise l'arbre B+.
\begin{enumerate}
   \item Les noeuds contiennent les valeurs de l'attribut (ou des
          attributs) cl\'e(s).
   \item Les feuilles contiennent chaque valeur index\'ee et le 
        $ROWID$ correspondant.
\end{enumerate}
Un index est stock\'e dans un segment qui lui est propre. On peut
le placer par exemple sur un autre disque que celui contenant 
la table.
\end{slide}
\begin{slide}
\subsection{Les $cluster$}

Le $cluster$ (regroupement) est un structure permettant d'optimiser
les jointures. Par exemple, pour les tables $CINEMA$ et $SALLE$
qui sont fr\'equemment jointes sur l'attribut $ID-Cinema$\,:
\begin{enumerate}
   \item On groupe les n-uplets de $CINEMA$ et de $SALLE$
        ayant m\^eme valeur pour l'attribut $ID-cinema$.
   \item On stocke ces groupes de n-uplets dans les pages
         d'un segment sp\'ecial de type $cluster$.
   \item On cr\'ee un index sur $ID-cinema$.
\end{enumerate}
\end{slide}
\begin{slide}

\myfigure{oracluster.eps}

\end{slide}
\begin{slide}
\subsection{Le hachage}

On d\'efinit un $hash\ cluster$ d\'ecrivant les 
caract\'eristiques physiques de la table\,:
\begin{verbatim}
CREATE CLUSTER hash-cin�ma (ID-cin�ma NUMBER(10))
   HASH IS ID-cin�ma HASHKEYS 1 000 SIZE 2K
\end{verbatim} 
$HASH\ IS$ (optionnel) sp\'ecifie la cl\'e \`a hacher. \\
$HASHKEYS$ est le nombre de valeurs de la cl\'e de hachage.\\
$SIZE$ est la taille d'un n-uplet.
ORACLE d\'etermine le nombre de pages allou\'ees, ainsi que la fonction
de hachage.
On fait r\'ef\'erence au $hash\ cluster$ en cr\'eant une table.
\end{slide}
% \end{slide}%
% LocalWords:  relation-exem.ps FNOM FADRESSE Abounayan Meudon Cima Preblocs DB
% LocalWords:  Gennevilliers Sarnaco PNOM COUTS COMDE Ivry Courbevoie QUANTITE
% LocalWords:  fourcli.ps Nuplet arit\'e d'arit\'e d'Attributs COUT CADRESSE to
% LocalWords:  FPCC maj l'Alg\`ebre unaires nuplet concat\`ene  width NUM S.A
% LocalWords:  R.A R.B nuplets Qu'est-ce peut-on Ingres \'Equijointure S.B FIND
% LocalWords:  d'\'Equijointure Durand Papin parpaings.ps PRODUIT.pnom ANY GET
% LocalWords:  FIRST WITHIN OWNER PRINT CLIENT.cnom CLIENT.cadresse NEXT em RBE
% LocalWords:  ensembliste Semijointure semijointure dupliqu\'es QTE PROD c.a.d
% LocalWords:  archi-optim.ps QBE n'appara qu ANNEE re division.ps




